\section{Method}

\subsection{Overview}

LEKAv2 replaces standard attention with a hybrid pipeline:

\begin{enumerate}
\item \textbf{Project} Q, K, V (frozen base + LoRA).
\item \textbf{Map} to dual-scale RFF features.
\item \textbf{Select} landmarks via a learned scorer.
\item \textbf{Compute} global landmark attention with \textbf{softpick}.
\item \textbf{Compute} local sliding-window attention.
\item \textbf{Fuse} via per-head gates.

\end{enumerate}
\textbf{Figure 4.1} shows the pipeline.

\begin{lstlisting}
    Input X $\in$ R^{B$$\times$L$$\times$d}
              |
         Q/K/V Projection (frozen + LoRA)
              |
       +------+------+
                      
      Q      K      V
       |      |      |
                     |
      (Q)    (K)     |  (Dual-scale RFF)
       |      |      |
       |             |
       |  Landmark  |
       |  Scorer    |
       |      |      |
       |             |
       |  K_land, V_land
       |      |      |
       +------+------+
               
      +---------------+
      |  Softpick     |
      |  Attention    |
      |  (Global)     |
      +-------+-------+
              |
      +-------+-------+
                       
    Local   Global   Gate
    Attn    Attn    ($\sigma$(w_h))
      |       |       |
      +-------+-------+
               
      (1-g)Local + g Global
              |
               
          out_proj
              |
               
           Output Y
\end{lstlisting}

\subsection{Dual-Scale Random Features}

Two projections:

\[
W_1 \in \mathbb{R}^{d_h \times F}, \quad W_2 \in \mathbb{R}^{d_h \times F/2}
\tag{4.1}
\]

Both initialized as $\mathcal{N}(0, 1/\sqrt{d_h})$ (scaled Gaussian).

\textbf{Feature map:}

\[
\phi(x) = \frac{1}{\sqrt{F}} [\cos(x W_1) \| \sin(x W_1)] \| \frac{1}{\sqrt{F/2}} [\cos(x W_2) \| \sin(x W_2)]
\tag{4.2}
\]

Output dimension: $2F + F = 3F$.

\textbf{Why two scales?} The two projections capture different frequency 
bands, similar to multi-resolution analysis.

\subsection{Landmark Distillation}

\textbf{Scorer.} A small MLP:

\[
s_\theta(\phi(k)) = \text{MLP}(\text{LN}(\phi(k))) \in \mathbb{R}
\tag{4.3}
\]

with architecture: LayerNorm($3F$) $\to$ Linear($3F$, 64) $\to$ GELU $\to$ Linear(64, 1).

\textbf{Selection.}

\[
L_\text{max} = \min\!\left(\max(L_\text{min}, \lceil \rho L \rceil), L\right)
\tag{4.4}
\]

\[
\mathcal{S}_\text{land} = \text{TopK}(s_\theta, L_\text{max})
\tag{4.5}
\]

\textbf{Position 0 guarantee.} We set $s_\theta[:, 0] \leftarrow 10^4$ to force 
retention.

\textbf{Gather.} $K_\text{land} = K[\mathcal{S}_\text{land}]$, $V_\text{land} = V[\mathcal{S}_\text{land}]$.

\subsection{Correction Matrix}

Per-head learnable correction applied to landmark values:

\[
V_\text{land}^\text{adj} = V_\text{land} \cdot W_\text{corr}^h
\tag{4.6}
\]

where $W_\text{corr} \in \mathbb{R}^{H \times d_h \times d_h}$ is 
initialized to the identity per head.

\textbf{Purpose.} Allows the model to learn head-specific value transformations 
during recovery.

\subsection{Softpick Attention}

\textbf{Scores.}

\[
\text{scores}_{ij} = \phi(q_i)^\top \phi(K_\text{land})_j^\top
\tag{4.7}
\]

\textbf{Masking.} Padding and causal masks applied.

\textbf{Softpick.}

\[
m_i = \max_j \text{scores}_{ij}, \quad
p_{ij} = \text{ReLU}\!\left(e^{(\text{scores}_{ij} - m_i)/T} - e^{-m_i/T}\right)
\tag{4.8}
\]

with $T = 1.0$ and fallback to uniform when $\sum_j p_{ij} < 10^{-6}$.

\textbf{Global context.}

\[
A_\text{global} = \sum_j p_{ij} V_{\text{land}, j}^\text{adj}
\tag{4.9}
\]

\subsection{Local Sliding Window}

\textbf{Window mask.}

\[
\text{mask}_{ij} = \begin{cases} 1 & |i - j| \leq W/2 \\ 0 & \text{otherwise} \end{cases}
\tag{4.10}
\]

\textbf{Local scores.} Standard scaled dot-product:

\[
\text{local}_{ij} = \frac{q_i \cdot k_j}{\sqrt{d_h}} \cdot \text{mask}_{ij}
\tag{4.11}
\]

Then causal + padding masks.

\textbf{Local attention.}

\[
A_\text{local} = \text{softmax}(\text{local}) \cdot V
\tag{4.12}
\]

\subsection{Local-Global Fusion}

Per-head gate:

\[
g_h = \sigma(w_h), \quad w_h^{(0)} = -2
\tag{4.13}
\]

\textbf{Fusion.}

\[
\text{out}_h = (1 - g_h) \cdot A_\text{local}^h + g_h \cdot A_\text{global}^h
\tag{4.14}
\]

With $g_h^{(0)} \approx 0.12$, the local branch dominates initially.

\subsection{Algorithm}

\textbf{Algorithm 1:} LEKAv2 forward pass
\begin{lstlisting}
Input:  X $\in$ R^{B$$\times$L$$\times$d}, attention_mask, position_embeddings
Output: Y $\in$ R^{B$$\times$L$$\times$d}

 1: Q $\leftarrow$ query(X);  K $\leftarrow$ key(X);  V $\leftarrow$ value(X)
 2: reshape to (B, H, L, d_h)
 3: if RoPE: Q, K $\leftarrow$ rotary(Q, K)
 4: if GQA:  K, V $\leftarrow$ repeat_kv
 5:  _Q $\leftarrow$ dual_rff(Q)                   (B, H, L, 3F)
 6:  _K $\leftarrow$ dual_rff(K)
 7: L_max $\leftarrow$ max(min_landmarks,   L)
 8: scores $\leftarrow$ landmark_scorer( _K)       (B, H, L)
 9: scores[:, :, 0] $\leftarrow$ 1e4
10: idx $\leftarrow$ topk(scores, L_max)
11: K_land $\leftarrow$ gather( _K, idx); V_land $\leftarrow$ gather(V, idx)
12: V_land_adj $\leftarrow$ V_land   W_corr
13: A_glob $\leftarrow$ softpick( _Q @ K_land^T) @ V_land_adj
14: A_loc $\leftarrow$ softmax(masked(QK^T)) @ V    window + causal
15: g $\leftarrow$ sigmoid(local_gate)              (H,)
16: out $\leftarrow$ (1-g) A_loc + g A_glob
17: reshape to (B, L, d); out_proj
18: return Y
\end{lstlisting}

\subsection{Complexity Analysis}

Per block:

\begin{table}[H]
\centering
\small
\begin{tabular}{lll}
\toprule
\textbf{Component} & \textbf{Time} & \textbf{Params} \\
\midrule
Q/K/V projections & $O(L d^2)$ & $3 d^2$ \\
Dual RFF & $O(L d (3F))$ & 0 (buffers) \\
Landmark scorer & $O(L \cdot 3F)$ & $O(F)$ \\
Landmark attention & $O(L \cdot L_\text{max} \cdot 3F)$ & 0 \\
Local window attention & $O(L \cdot W \cdot d_h)$ & 0 \\
Correction matrix & $O(L_\text{max} \cdot d_h^2)$ & $H d_h^2$ \\
Fusion & $O(L d)$ & $H$ \\
Output & $O(L d^2)$ & $d^2$ \\
\textbf{Total} & $O(L d^2 + L \cdot L_\text{max} \cdot F)$ & $\sim 4 d^2 + H d_h^2$ \\
\bottomrule
\end{tabular}
\end{table}

For $\rho = 0.5$ and $L \gg F$, the landmark term is 
$O(L^2 \cdot F)$ -- \textbf{quadratic} in $L$! But with $L_\text{max} = \rho L$ 
and $\rho = 0.5$, we get $O(L^2)$ again.

\textbf{Note.} LEKAv2 is \textbf{not truly linear} unless we set $L_\text{max} = $ 
constant. The current implementation uses $L_\text{max} = \rho L$ -- this 
is effectively \textbf{linear-quadratic} (still $O(L^2)$ but with a smaller 
constant).

\subsection{Parameter Analysis}

For $d = 768$, $H = 12$, $d_h = 64$, $F = 128$, LoRA $r_L = 16$:

\begin{table}[H]
\centering
\small
\begin{tabular}{lll}
\toprule
\textbf{Component} & \textbf{Baseline} & \textbf{LEKAv2} \\
\midrule
Q/K/V/O & $4 d^2 = 2.36M$ & $2.36M$ (frozen) \\
LoRA & 0 & $98K$ \\
RFF buffers & 0 & 0 (not params) \\
Landmark scorer & 0 & ~$9K$ \\
Correction matrix & 0 & $H d_h^2 = 49K$ \\
Local gate & 0 & 12 \\
\textbf{Total per layer} & 2.36M & 2.51M \\
\textbf{Increase} & -- & \textbf{+6\%} \\
\bottomrule
\end{tabular}
\end{table}

Modest parameter increase.

\subsection{Training}

Trainable:
\begin{itemize}
\item LoRA adapters (frozen base).
\item Landmark scorer.
\item Correction matrix.
\item Local gate.

\end{itemize}
Frozen:
\begin{itemize}
\item Q/K/V/O base weights.
\item Random projection buffers (not parameters).
\end{itemize}
\begin{lstlisting}

---
\end{lstlisting}